\section{Dataset and Exploratory Analysis}
\label{sec:dataset}

Data were collected from public CompetitionSuite recap pages (a web platform through which judges submit scores at SCPA events) using
the cs-parse library~\cite{scpadb2026, csparse2026}. The resulting database
contains 4,441 performance records and 56,910 score records spanning the
2017--2026 SCPA seasons, excluding 2021, when no season was held. Each
performance record represents one ensemble's scored appearance at a
competition. Music, Visual, Music Effect, and Visual Effect are judged
separately; championship finals may use multiple adjudicators for each caption.
Ensembles compete within a \emph{class}: a
competitive division grouping programs by skill level. Independent ensembles
compete in A (basic concepts and skills), Open (intermediate), or World
(advanced) classes (coded PIA, PIO, PIW); scholastic ensembles, which are
school-affiliated programs, compete in parallel scholastic classes (PSA, PSO,
PSW) or the Junior division for middle-school programs (PSJ, also evaluated on
basic concepts and skills similar to A class).

Three concert classes (PSCA, PSCO, PSCW) use a stationary, non-marching format. They are excluded because their caption structure omits the marching
visual framework and is incompatible with the four-caption marching
sheet~\cite{wgi2026scoresheets}. Sixteen
all-zero records are retained in the public database for historical
completeness but excluded from all analysis; thirteen are marching records. The
seven marching classes contain 3,808 raw performance records and 52,288 score records;
after removing the thirteen marching artifacts, the clean analysis scope
contains 3,795 performance records. The 2020 season, curtailed before its
championship, is included in descriptive analysis only. The 2026 season is
complete through April 11, 2026.

An \emph{ensemble-season} is one performing ensemble observed across one
season. Its \emph{debut} is its first valid, non-championship scored appearance
within that season. The primary modeling target is the terminal
championship subtotal score (the sum of adjudicated caption scores at the
furthest championship stage the ensemble reached), excluding timing penalties
(deductions for timing, logistics, or rule compliance rather than artistic merit)
and bounded to $[0, 100]$. Championship competition may include prelims,
semifinals, and finals, so ``terminal'' refers to the latest stage reached. The
modeling cohort requires each ensemble-season to have both a debut and a
championship result, and excludes ensembles that
appeared only at championships (typically visiting groups without a
regional-season debut). An ensemble-season is also excluded if that same
ensemble competes in multiple classes during the season, because its terminal
class is not known at debut. A program may still contribute separate
ensemble-seasons when it fields distinct ensembles concurrently (for example, a school entering both a Scholastic A and a Scholastic Open ensemble in the same season). After applying
these rules and a requirement of at least two ensemble-seasons per class-season (needed for the percentile formula denominator), 781 ensemble-seasons remain:
574 development observations from 2017--2019 and 2022--2024, and 207 held-out
observations from 2025--2026. Only ensemble-seasons with a championship result
receive a target, which induces a selection effect the model cannot remove.

Descriptive scores rise strongly within most completed seasons: the median
subtotal increases from roughly 70--74 points in the first competition week to
approximately 87--90 points near championships. Terminal championship medians
also differ substantially by class and year, motivating the use of class,
timing, and competition-context features rather than treating a debut score as
directly comparable across all observations.
Figure~\ref{fig:class-score-trends} shows median terminal championship
subtotals by class across seasons.

% Median terminal championship subtotal by class and season.
% Data: ../analysis/outputs/tables/terminal_championship_medians.csv
% Wide-format CSV; one \addplot per marching class (concert classes excluded).
\begin{figure}[t]
\centering
\pgfplotstableread[col sep=comma]{../analysis/outputs/tables/terminal_championship_medians.csv}\medians
\begin{tikzpicture}
\begin{axis}[
  width=\columnwidth,
  height=5.0cm,
  ylabel={Median subtotal},
  xmin=2016.5, xmax=2026.5,
  xtick={2017,2018,2019,2022,2023,2024,2025,2026},
  xticklabels={2017,2018,2019,2022,2023,2024,2025,2026},
  x tick label style={rotate=45, anchor=east, font=\small},
  ymin=74, ymax=97,
  ytick={75,80,85,90,95},
  ylabel style={font=\small},
  tick label style={font=\small},
  legend style={
    font=\footnotesize,
    at={(0.5,-0.16)},
    anchor=north,
    legend columns=4,
    cells={anchor=west},
    column sep=4pt,
    draw=none,
    fill=none,
  },
  grid=major,
  grid style={dotted, gray!40},
  cycle list name=color list,
]
\addplot+[mark=*, mark size=1.4pt] table[x=season_year, y=piw]{\medians}; \addlegendentry{PIW}
\addplot+[mark=square*, mark size=1.3pt] table[x=season_year, y=psw]{\medians}; \addlegendentry{PSW}
\addplot+[mark=triangle*, mark size=1.4pt] table[x=season_year, y=pio]{\medians}; \addlegendentry{PIO}
\addplot+[mark=diamond*, mark size=1.4pt] table[x=season_year, y=pso]{\medians}; \addlegendentry{PSO}
\addplot+[mark=pentagon*, mark size=1.3pt] table[x=season_year, y=psj]{\medians}; \addlegendentry{PSJ}
\addplot+[mark=*, mark size=1.4pt, dashed] table[x=season_year, y=pia]{\medians}; \addlegendentry{PIA}
\addplot+[mark=square*, mark size=1.3pt, dashed] table[x=season_year, y=psa]{\medians}; \addlegendentry{PSA}
\end{axis}
\end{tikzpicture}
\caption{Median terminal championship subtotal by class, 2017--2026 (2020 and 2021 excluded because neither has a championship result). Medians vary by class and season, supporting class-aware modeling.}
\Description{Line chart of median terminal championship subtotal by class and season. Class medians vary across years, with World classes generally near the upper end of the plotted range.}
\label{fig:class-score-trends}
\end{figure}

